\section{A conditional tape-context fold from 209 to 208 states}
\label{sec:local208}

The reduction in this section acts directly on the supplied Joint U209
transition table. It preserves the tape alphabet, the blank initializer, and
the entire numerical source program. Its additional information is a property
of the recent computation: two states with different zero branches can share
one physical state because their reachable tape contexts distinguish those
branches. The required property of one state is conditional on its scanned
symbol. The opposite-predecessor-bit interfaces considered in the supplied
searches do not include this case.

All state labels in this section are \emph{native} labels. In particular, 211
is a working state and 254 is the separate halt label. Dense renumbering is
performed only after the reduction. Let $\Uold$ denote the Joint parent and
$\Unew$ the reduced machine.

\subsection{The complete modification}
\label{subsec:local208-patch}

Delete state 172 and redirect its sole incoming transition to state 211.
Change the zero transition of state 211 as follows:
\begin{center}
\begin{tabular}{llll}
\toprule
Row & Parent action & Reduced action & Purpose\\
\midrule
$116/0$ & $0\,R\,172$ & $0\,R\,211$ & Redirect the deleted state\\
$211/0$ & $0\,R\,157$ & $0\,L\,217$ & Inspect the preceding tape context\\
\bottomrule
\end{tabular}
\end{center}
The removed rows are $172/0:0\,L\,67$ and $172/1:1\,R\,184$.
Every other retained row is unchanged. Thus 209 working states and 418
transition rows become 208 working states and 416 rows. The initial state
remains 0, and no halt transition is added or removed.

For the simulation, define the state map
\[
 f(q)=\begin{cases}211,&q=172,\\q,&q\ne172.\end{cases}
\]
The map identifies two states, but changes neither the head position nor
the tape. Its lack of injectivity is harmless: the proof follows a definite
parent computation and establishes agreement at explicitly defined times.

\subsection{Reachable contexts, proved by induction}
\label{subsec:local208-guards}

A configuration is $(q,h,T)$, where $h\in\mathbb Z$ and
$T:\mathbb Z\to\{0,1\}$. Write $t_j=T(h+j)$ for the tape viewed from
the current head. The following five conditions are imposed only when the
configuration has the indicated state:
\begin{align*}
 I_{116}:&\quad t_{-1}=1, &
 I_{172}:&\quad t_{-2}=1\ \land\ t_{-1}=0,\\
 I_{217}:&\quad t_{-1}=1, &
 I_{163}:&\quad t_{-1}=1\ \lor\ t_{-2}=1,\\
 I_{211}:&\quad t_0=1\ \lor\ t_{-1}=1.
\end{align*}
Let $I(q,h,T)$ be the conjunction of the five corresponding implications
from the state equality. In particular, $I_{211}$ says that a scanned zero
at state 211 forces the cell immediately to its left to be one. It does
\emph{not} assert that every entry to 211 has a left-hand one.

The entire incoming-edge census needed for these conditions is small:
\begin{center}
\begin{tabular}{cl}
\toprule
Target & Every parent transition entering the target\\
\midrule
172 & $116/0:0\,R\,172$\\
116 & $215/1:1\,R\,116$\\
211 & $163/0:0\,L\,211$, \quad $163/1:1\,R\,211$\\
163 & $132/1:1\,R\,163$, \quad $199/1:1\,R\,163$,
      \quad $217/0:0\,R\,163$\\
217 & $216/1:1\,R\,217$\\
\bottomrule
\end{tabular}
\end{center}
These are inventories of all physical transition rows, including entries
arising outside ordinary source decoding. The argument therefore does not
depend on a sampled run or an assumption that additional entrances are
unreachable.

\begin{lemma}[Five-state invariant]
\label{lem:local208-invariant}
Every configuration of $\Uold$ reached from state 0 satisfies $I$, on
every initial binary tape.
\end{lemma}
\begin{proof}
The initial configuration satisfies $I$ because 0 is none of the five
listed states. Assume $I$ holds before a transition. Conditions whose
target state is not reached impose no obligation afterward, so it suffices
to inspect the displayed census.

An entry to 116 or 217 writes one and moves right, establishing its
left-hand-one condition immediately. An entry to 172 comes from 116/0.
It writes zero immediately left of the new head; the previous left-hand
one supplied by $I_{116}$ becomes the cell two positions left of that head.
This proves $I_{172}$.

An entry to 163 from 132/1 or 199/1 writes one immediately left of the new
head. Its remaining entrance, 217/0, writes a left-hand zero but carries
the one supplied by $I_{217}$ to the position two cells left. Each case
establishes the disjunction $I_{163}$.

Finally, entry to 211 through 163/1 writes one immediately left of the
new head. Entry through 163/0 moves left instead. The new current cell is
the old $t_{-1}$, and the new left cell is the old $t_{-2}$. Their disjunction
is exactly $I_{163}$, so $I_{211}$ follows in this case as well. This proves
preservation of all five conditions and completes the induction.
\end{proof}

The invariant explains the available distinction. State 172 necessarily
arrives immediately to the right of the word $10$. State 211, whenever it
reads zero, instead has a one immediately to its left. Recovering this
information needs at most three preceding parent transitions. Nevertheless,
the inductive formulation is preferable for composition: it applies uniformly
at every time, including the beginning of the run, without a separate
assumption that a sufficiently long history already exists.

\subsection{Four exact finite identities}
\label{subsec:local208-blocks}

At a simulation boundary the reduced machine has state $f(q)$, head $h$,
and the same tape $T$ as the parent. Only parent states 172 and 211 require
special treatment. The four contracts below specify the complete effects
of their replacement blocks. Head displacements are relative to $h$.
\begin{center}
\begin{tabular}{ccccc}
\toprule
Parent state & Read & Additional tape condition & Steps & Final state, head\\
\midrule
172 & 0 & $t_{-2}=1,\ t_{-1}=0$ & 7 & $67,\ -1$\\
172 & 1 & None & 1 & $184,\ +1$\\
211 & 0 & $t_{-1}=1$ & 3 & $157,\ +1$\\
211 & 1 & None & 1 & $184,\ +1$\\
\bottomrule
\end{tabular}
\end{center}
Every contract preserves the entire tape. The additional conditions for
the zero cases follow from Lemma~\ref{lem:local208-invariant}. The one
cases need only the stated scanned bit, even though other invariant
conditions may also hold.

\begin{lemma}[Local replacement]
\label{lem:local208-replacement}
Suppose a nonhalting parent configuration $C=(q,h,T)$ satisfies $I$,
and its successor is $C'=(q',h',T')$. From $(f(q),h,T)$ the reduced machine
reaches $(f(q'),h',T')$ in one, three, or seven transitions. All earlier
configurations of this block are nonhalting.
\end{lemma}
\begin{proof}
If $q\notin\{172,211\}$, the parent and reduced rows agree after applying
$f$ to the target. This includes row 116/0, whose target is precisely the
one redirected by $f$. There are 414 such ordinary rows.

For either exceptional state reading one, the reduced machine uses the
unchanged row $211/1:1\,R\,184$. Both parent states have that same row,
so one transition suffices.

For parent state 211 reading zero, its invariant supplies $t_{-1}=1$.
The new row moves left into 217; row $217/1:1\,R\,215$ moves back to
the original zero; row $215/0:0\,R\,157$ gives the endpoint of the
parent's $211/0$ transition. The three transitions read $0,1,0$ and write
those same bits.

For parent state 172 reading zero, the invariant supplies
$(t_{-2},t_{-1},t_0)=(1,0,0)$. The full reduced computation is:
\begin{center}
\begin{tabular}{rrrrrr}
\toprule
Step & State before & Head before & Read/write & State after & Head after\\
\midrule
1 & 211 & $0$  & $0/0$ & 217 & $-1$\\
2 & 217 & $-1$ & $0/0$ & 163 & $0$\\
3 & 163 & $0$  & $0/0$ & 211 & $-1$\\
4 & 211 & $-1$ & $0/0$ & 217 & $-2$\\
5 & 217 & $-2$ & $1/1$ & 215 & $-1$\\
6 & 215 & $-1$ & $0/0$ & 157 & $0$\\
7 & 157 & $0$  & $0/0$ & 67  & $-1$\\
\bottomrule
\end{tabular}
\end{center}
Its endpoint is exactly the parent transition $172/0:0\,L\,67$.
Every read is of one of the three specified cells and every write preserves
that cell. The three-step case has the same property for its two cells.
No unspecified exterior cell is inspected, and no exceptional block visits
halt label 254. This proves all cases.
\end{proof}

The seven-step block visits the merged state 211 twice. There is no
circular proof here: all seven transitions have been evaluated explicitly.
In particular, the argument does not assume that arbitrary internal
configurations of the reduced computation satisfy the parent invariant.
The invariant is required only at parent simulation boundaries.

\subsection{Global equivalence and an exact clock}
\label{subsec:local208-global}

For a parent computation, let $N_{172,0}(n)$ count occurrences of state 172
reading zero among its first $n$ transitions, and define $N_{211,0}(n)$
similarly. Counts are taken along the parent run, not the internal states
of the reduced blocks. Define
\begin{equation}
\label{eq:local208-clock}
 \tau(n)=n+6N_{172,0}(n)+2N_{211,0}(n).
\end{equation}
The excesses six and two are the extra transitions used by the seven-step
and three-step replacements.

\begin{theorem}[Uniform halting equivalence]
\label{thm:local208-equivalence}
For every binary tape $T$ and initial head position $h$, $\Uold$ and
$\Unew$, both started in state 0 on $(h,T)$, either both halt or both
run forever. At every parent time $n$ up to and including its first halt,
the reduced configuration at time $\tau(n)$ has the same head and tape,
and its state is the image under $f$ of the parent state. Furthermore,
\[
 n\le\tau(n)\le7n.
\]
If the parent first halts at time $H$, the reduced machine first halts
at exactly $\tau(H)$.
\end{theorem}
\begin{proof}
The configurations agree initially because $f(0)=0$. Induct on the number
of parent transitions. Lemma~\ref{lem:local208-invariant} supplies the
required invariant at each parent configuration, and
Lemma~\ref{lem:local208-replacement} extends the agreement by the
appropriate positive block. The block lengths give
\eqref{eq:local208-clock}. Each length is between one and seven, yielding
the bounds. No block contains an earlier halt, so a parent first halt
produces a reduced first halt exactly at the corresponding boundary.

If the parent never halts, the same induction supplies nonhalting reduced
blocks at all boundaries. Their strictly positive lengths make
$\tau(n)$ unbounded, and together the blocks cover every reduced time.
The reduced machine therefore cannot halt between boundaries or after
all of them. This also proves the converse implication of the halting
equivalence. Determinism ensures that these blocks are prefixes of the
actual reduced run, rather than merely possible executions.
\end{proof}

Thus the result concerns the whole computation, not only a finite prefix.
It also supplies quantitative finite statements: a parent halt within $n$
steps implies a reduced halt within $7n$ steps; a reduced halt within $n$
steps implies a parent halt within $n$ steps. No consistency assumption
appears anywhere in the reduction.

\subsection{Arithmetic scope and the other U209 family}
\label{subsec:local208-scope}

The arbitrary-tape theorem above is a mathematical statement about tape
functions. For the required blank input, configurations and finite runs
have the usual arithmetic encodings. The five invariant checks, the four
finite identities, and the induction on coded computation length are
formalizable in $\PA$. The explicit bound $7n$ removes any need for an
unbounded search to locate a corresponding finite prefix. Consequently,
this construction supplies the PA-formalizable equivalence
\[
 \PA\vdash
 \bigl(\Halt(\Uold,\mathrm{blank})
       \leftrightarrow\Halt(\Unew,\mathrm{blank})\bigr).
\]
Composing it with a justified parent consistency calibration transfers
that calibration to $\Unew$. The reduction neither establishes
$\Con(\PA)$ nor replaces the separate proof-search and theory
interpretation obligations. The delivered mathematical proof and finite
certificates do not constitute a newly exported proof-assistant certificate
of the complete consistency endpoint.

The distinct Compatibility U209 family admits the same reduction:
delete 172 into 211 and replace row $211/0$ by $0\,L\,214$.
Its local roles correspond as follows:
\begin{center}
\begin{tabular}{lcc}
\toprule
Role & Joint parent & Compatibility parent\\
\midrule
Left-context selector & 217 & 214\\
State feeding the merged-state entries & 163 & 216\\
Right branch of the selector & 215 & 163\\
Predecessor writing one before the selector & 216 & 217\\
\bottomrule
\end{tabular}
\end{center}
States 116, 172, 211, 157, 67, and 184 retain their roles. The corresponding
complete censuses are: 172 from 116/0; 116 from 163/1; 211 from 216/0
and 216/1; 216 from 132/1, 199/1, and 214/0; and 214 from 217/1.
Their directions and written bits are exactly those used in the invariant
proof. The seven-step state sequence becomes
\[
 211\to214\to216\to211\to214\to163\to157\to67.
\]
The same tape conditions, head positions, block lengths, and clock formula
therefore apply. This produces a second 208-state table from the other
supplied machine family, with no additional source transformation.
